\subsection{Epistemic noninterference on the composed v3 authority surface}

The v3 experience substrate is designed to change future search behaviour without allowing experience-side state to silently become scientific authority.  The executable specification names this property \emph{epistemic noninterference}.  Let $\pi_{\mathrm{auth}}(X)$ denote the scientific-authority projection over registered grounding, representation, mechanism, identification and decision-use coordinates.  For a transition sequence composed only of experience, retrieval, workspace, strategy, reflection, routing or training-allocation operations, the intended invariant is
\[
\pi_{\mathrm{auth}}(\tau_E^*(X_t))=\pi_{\mathrm{auth}}(X_t),
\]
unless the sequence contains a separately registered evidence-bearing scientific-authority transition.  The projection is non-monotone: a properly attested refutation may revoke active authority while preserving the historical certificate and challenge record.

The historical pre-integration runtime returned \path{NO_INTEGRATION_SURFACE} because experience state and scientific authority were not yet composed.  That result is preserved as negative/incomplete history rather than retrospectively relabelled as a pass.  The current revision supersedes that boundary.  \texttt{RAKLV3State} now carries an immutable \texttt{ScientificAuthorityProjection} containing canonical claim atoms, content-bound scientific-evidence bindings, certificate/event history and the active scientific-authority view.  The integration is deliberately immutable so a transition cannot mutate a shared ledger behind both the pre- and post-transition projections.

Scientific authority is certificate-bound rather than declaration-bound.  Promotion resolves a protected attestation whose subject binds the actual claim text, requested axis, scope and registered evidence digests.  Method-authority attestations for Lessons or ResearchTools are not accepted as scientific-authority attestations.  Public v3 transitions are classified by ownership, and only the registered scientific promotion, revocation and supersession transitions own the authority coordinate.  Registering a claim or evidence object changes canonical scientific content but does not by itself grant authority.

Historical benchmark identity is also protected.  The original state fingerprint remains the v1 experience/framework fingerprint, while \texttt{state\_fingerprint\_v2} explicitly includes the scientific-authority projection.  Authority-sensitive experiments must therefore pin v2 instead of silently changing the byte identity of older benchmark packets.

\paragraph{Additive v3 commitment migration.}
A third, explicitly versioned commitment domain is reserved for dual-write migration and cross-runtime replay. The v3 commitment surface digests named components---each bound to its own schema version---into a sorted, unique, hash-chained state commitment that carries a sequence number and the digest of its predecessor, so migration is forward-linkable rather than a silent re-keying of historical \texttt{state\_fingerprint} or \texttt{state\_fingerprint\_v2} values. The commitment profile preserves exact Unicode scalar sequences: any field that semantically requires NFC normalization must enforce that invariant at its own schema boundary, so an integrity commitment neither rewrites nor globally forbids older state. Like v2 before it, the v3 domain is an additive identity layer; it grants no scientific authority and exists so that authority-sensitive experiments can pin an exact commitment version without retroactively altering the byte identity of prior benchmark packets.

\begin{sloppypar}
The executable checker in \path{src/rakl/epistemic_noninterference.py}, together with \path{src/rakl/v3_scientific_authority.py}, now exercises the invariant against this real composition surface.  Planted hostile worlds cover experience-to-evidence leakage, repetition/routing/reflection/workspace/self-evolution shortcuts, failure-to-impossibility, provenance-to-independence, prediction-to-mechanism, mechanism-to-identification and unattested authority movement.  Benign controls require both that non-authority state actually changes and that a legitimate evidence-bearing promotion or refutation can still move the authority projection.  Integration-path mutation tests separately plant bypasses in the authority contract and require the suite to detect them.
\end{sloppypar}

This result is still narrower than a behavioural deployment claim.  It establishes that the composed runtime has an executable protected authority interface under the registered transitions.  It does not establish that a language model will rarely propose unsupported claims, that Orion has lower behavioural authority leakage than strong controls, or that any particular base model is capable of executing the epistemic protocol.  Those are separately measured questions.

The same boundary extends to the optional training projection introduced in Section~\ref{sec:training-time-projection}.  A model checkpoint, mastery estimate or gradient-allocation proposal may change dramatically while the scientific-authority projection remains unchanged.  A training run may later provide evidence for a claim \emph{about model behaviour} only through a separately registered evaluator/run/artifact binding and the ordinary scientific promotion path.  Thus weight updates, like experience and retrieval, are behaviour-changing operations rather than scientific-authority shortcuts.

V3 additionally separates retained experience from canonical admission.  A \texttt{TaskEpisode} may be preserved as \path{PROPOSAL_SHADOW_STORED} without satisfying the stronger \path{CANONICAL_INVENTORY_ADMITTED} gate, and protected consumers fail closed if a shadow object is presented as canonical inventory.  Persistence, retrieval, training exposure and repeated success therefore remain distinct from scientific evidence and authority.
